\documentclass[11pt]{article}
\usepackage[a4paper,margin=1in]{geometry}
\usepackage{amsmath,amssymb,amsfonts,bm}
\usepackage{hyperref}
\usepackage{enumitem}
\usepackage{mathtools}

\title{Regulated Causal Defects as Alternatives to Singular Black-Hole Endpoints:\\A First Effective-Tensor Framework}
\author{First Draft}
\date{\today}

\begin{document}
\maketitle

\begin{abstract}
Classical general relativity predicts trapped regions and, under suitable energy and causality conditions, singularity formation during sufficiently compact gravitational collapse. This draft develops a speculative effective framework in which the black-hole endpoint is reinterpreted not as literal infinite-density spacetime failure, but as a regulated causal defect: a phase of spacetime in which extreme compactness activates additional stress-energy contributions that prevent curvature invariants from diverging. The proposal is deliberately conservative in its first formulation: the exterior geometry may approximate Schwarzschild or Kerr at accessible radii, while deviations appear only near an ultra-compact core or would be detectable through horizon-scale dynamics. We introduce a regulator tensor $T^{\rm reg}_{\mu\nu}$, define activation by a compactness/curvature threshold, derive minimal consistency conditions from covariant conservation, and list phenomenological signatures. The model is not presented as a completed theory of quantum gravity, but as a clean mathematical container for the hypothesis that collapse enters a ``cooldown deadlock'' phase rather than a singular endpoint.
\end{abstract}

\section{Motivation}
The standard semiclassical picture of gravitational collapse contains two sharply different regimes. Far from the collapsing body, the geometry is well approximated by classical general relativity. Near the would-be singular region, however, the extrapolation of classical equations becomes physically suspect because curvature invariants can diverge and quantum gravitational effects are expected to become important.

This draft formalizes the following hypothesis:
\begin{quote}
A black hole is not fundamentally a pointlike infinite-density object. It is an extreme causal phase of spacetime, triggered by compact energy concentration, in which a regulator sector stores, redistributes, or slowly leaks energy so that curvature remains finite.
\end{quote}

The goal is not to deny the empirical success of black-hole exterior solutions. Instead, the claim is that classical general relativity may correctly identify a collapse threshold while misidentifying the final internal phase.

\section{Background and related ideas}
This proposal overlaps with several existing lines of thought. Regular black holes replace the singularity by finite-curvature interiors. Gravastars replace the classical black-hole interior by a vacuum-like condensate phase \cite{MazurMottola2001}. Planck-star scenarios suggest that quantum gravity can halt collapse at high density and produce long-lived compact remnants \cite{RovelliVidotto2024}. More recent work also studies regular black-hole and gravastar-like solutions derived from modified-gravity or effective field-theory actions \cite{EichhornFernandes2025,Jampolski2025}.

The present draft differs mainly in language and organization. It frames the endpoint as a \emph{regulated causal defect}, using an effective stress-energy tensor as the formal object that turns the intuition into equations.

\section{Classical collapse threshold}
For a spherical mass distribution with total mass $M$ and areal radius $R$, define the compactness
\begin{equation}
    C = \frac{2GM}{Rc^2}.
\end{equation}
Classically, $C \simeq 1$ marks the Schwarzschild radius threshold. The classical Einstein equation is
\begin{equation}
    G_{\mu\nu} + \Lambda g_{\mu\nu}
    = \frac{8\pi G}{c^4}T_{\mu\nu}.
\end{equation}
The singularity theorems do not say that a singularity is directly observed. They show geodesic incompleteness under assumptions such as trapped surfaces, causal conditions, and energy conditions. The present framework targets those assumptions by adding an effective regulator sector at high compactness or high curvature.

\section{Effective regulated field equation}
We replace the classical source by an effective source:
\begin{equation}
    G_{\mu\nu} + \Lambda g_{\mu\nu}
    = \frac{8\pi G}{c^4}
    \left(
        T^{\rm m}_{\mu\nu}
        + T^{\rm q}_{\mu\nu}
        + T^{\rm reg}_{\mu\nu}
    \right),
    \label{eq:modified-einstein}
\end{equation}
where
\begin{itemize}[leftmargin=2em]
    \item $T^{\rm m}_{\mu\nu}$ is ordinary matter and radiation,
    \item $T^{\rm q}_{\mu\nu}$ is the semiclassical quantum-field contribution,
    \item $T^{\rm reg}_{\mu\nu}$ is the proposed regulator stress-energy.
\end{itemize}

The Bianchi identity requires
\begin{equation}
    \nabla^\mu
    \left(
        T^{\rm m}_{\mu\nu}
        + T^{\rm q}_{\mu\nu}
        + T^{\rm reg}_{\mu\nu}
    \right)=0.
    \label{eq:conservation}
\end{equation}
Therefore the regulator cannot be inserted arbitrarily. It must either be separately conserved,
\begin{equation}
    \nabla^\mu T^{\rm reg}_{\mu\nu}=0,
\end{equation}
or exchange energy-momentum with matter and quantum sectors:
\begin{equation}
    \nabla^\mu T^{\rm reg}_{\mu\nu}=-Q_\nu,
    \qquad
    \nabla^\mu(T^{\rm m}_{\mu\nu}+T^{\rm q}_{\mu\nu})=Q_\nu.
\end{equation}

\section{Activation rule}
Let $\mathcal I[g]$ be a scalar curvature diagnostic, for example
\begin{equation}
    \mathcal I[g]
    = \ell_*^4 R_{\alpha\beta\gamma\delta}R^{\alpha\beta\gamma\delta},
\end{equation}
where $\ell_*$ is the regulator length scale. A smooth activation function is
\begin{equation}
    A(\mathcal I)=\frac{1}{2}\left[1+\tanh\left(\frac{\mathcal I-\mathcal I_c}{\Delta}\right)\right].
\end{equation}
Then
\begin{equation}
    T^{\rm reg}_{\mu\nu}=A(\mathcal I)\,\Theta_{\mu\nu},
\end{equation}
where $\Theta_{\mu\nu}$ is the regulator tensor structure. The constants $\mathcal I_c$ and $\Delta$ define the onset and sharpness of the transition.

An alternative activation variable is compactness:
\begin{equation}
    A(C)=\frac{1}{2}\left[1+\tanh\left(\frac{C-C_c}{\delta C}\right)\right].
\end{equation}
Curvature activation is more covariant; compactness activation is simpler for spherical collapse.

\section{Minimal spherical model}
Use the static spherically symmetric ansatz
\begin{equation}
    ds^2=-e^{2\Phi(r)}c^2dt^2+\left(1-\frac{2Gm(r)}{c^2r}\right)^{-1}dr^2+r^2d\Omega^2.
\end{equation}
Let the effective density and radial pressure be
\begin{equation}
    \rho_{\rm eff}=\rho_m+\rho_q+\rho_{\rm reg},
    \qquad
    p_{r,\rm eff}=p_{r,m}+p_{r,q}+p_{r,\rm reg}.
\end{equation}
The mass function satisfies
\begin{equation}
    \frac{dm}{dr}=4\pi r^2\rho_{\rm eff}(r).
\end{equation}
The generalized Tolman--Oppenheimer--Volkoff equation becomes
\begin{equation}
    \frac{dp_{r,\rm eff}}{dr}
    = -\frac{(\rho_{\rm eff}c^2+p_{r,\rm eff})
    \left[Gm(r)+4\pi G r^3p_{r,\rm eff}/c^2\right]}
    {c^2r^2\left(1-2Gm(r)/(c^2r)\right)}
    + \frac{2}{r}\left(p_{t,\rm eff}-p_{r,\rm eff}\right).
\end{equation}
A finite-curvature endpoint requires
\begin{equation}
    \rho_{\rm eff}(0)<\infty,
    \qquad
    p_{r,\rm eff}(0)<\infty,
    \qquad
    p_{t,\rm eff}(0)<\infty
\end{equation}
with regularity condition
\begin{equation}
    m(r)\sim \frac{4\pi}{3}\rho_{\rm eff}(0)r^3
    \quad \text{as } r\to 0.
\end{equation}

\section{Cooldown-deadlock interpretation}
The phrase ``cooldown deadlock'' can be translated into a dynamical energy-balance equation. Let $E_{\rm core}$ be energy stored in the regulated core. Then
\begin{equation}
    \frac{dE_{\rm core}}{dt}
    = \dot E_{\rm in}-\dot E_{\rm leak}-\dot E_{\rm grav},
\end{equation}
where $\dot E_{\rm in}$ is infalling flux, $\dot E_{\rm leak}$ is outgoing quantum or nonclassical flux, and $\dot E_{\rm grav}$ is gravitational-wave or geometric dissipation.

A deadlock-like state is a metastable fixed point:
\begin{equation}
    \frac{dE_{\rm core}}{dt}\approx 0,
    \qquad
    E_{\rm core}\gg \dot E_{\rm leak}\tau_{\rm dyn},
\end{equation}
where $\tau_{\rm dyn}$ is the collapse timescale. This differs from ordinary Hawking evaporation unless the leakage channel is stronger than the standard Hawking flux.

\section{Phenomenological signatures}
The model should be constrained by at least four observational channels:
\begin{enumerate}[leftmargin=2em]
    \item \textbf{Exterior metric tests.} Far from the compact object, deviations from Kerr must be small.
    \item \textbf{Ringdown.} Modified interiors or near-horizon structure may shift quasinormal modes or produce late-time echoes.
    \item \textbf{Accretion images.} Horizon-scale emission and shadow structure must remain consistent with Event Horizon Telescope constraints.
    \item \textbf{Evaporation or leakage.} If the model predicts enhanced leakage, it must not contradict astrophysical black-hole longevity.
\end{enumerate}

\section{Falsifiability}
This framework fails if every mathematically consistent $T^{\rm reg}_{\mu\nu}$ either violates known constraints strongly, produces exterior deviations already ruled out, or reduces to a relabeling of already known regular black-hole models without new predictions.

\section{Conclusion}
The regulated-causal-defect model is a formal container for the intuition that black holes are not literal singular endpoints but extreme causal phases. Its core mathematical object is the regulator tensor $T^{\rm reg}_{\mu\nu}$, activated by compactness or curvature. The next step is to choose an action principle that generates $T^{\rm reg}_{\mu\nu}$ rather than inserting it phenomenologically.

\begin{thebibliography}{9}
\bibitem{MazurMottola2001}
P. O. Mazur and E. Mottola,
``Gravitational Condensate Stars: An Alternative to Black Holes,''
\texttt{arXiv:gr-qc/0109035}.

\bibitem{RovelliVidotto2024}
C. Rovelli and F. Vidotto,
``Planck stars, White Holes, Remnants and Planck-mass quasi-particles,''
\texttt{arXiv:2407.09584}.

\bibitem{EichhornFernandes2025}
A. Eichhorn and P. G. S. Fernandes,
``Regular black holes without mass-inflation instability and gravastars from modified gravity,''
\texttt{arXiv:2508.00686}.

\bibitem{Jampolski2025}
D. Jampolski,
``On the formation of gravastars,''
\texttt{arXiv:2509.15302}.

\bibitem{Visser2009}
M. Visser,
``Black holes in general relativity,''
\texttt{arXiv:0901.4365}.
\end{thebibliography}

\end{document}
